\section{Worked problems: napkin math at three scales}

\subsection{Llama 3 70B on eight H100s}
\textbf{Problem.} Llama~3~70B (80 layers, width 8,192, 64 query and 8 key/value
heads of 128, FFN 28,672, vocabulary 128,256, untied tables) runs in FP8 ---
weights and cache --- on eight H100s joined by NVLink, split by tensor
parallelism. Ignoring communication, predict the step time, per-user rate,
aggregate throughput and cost at batch one and at the largest batch that fits at
4,096 tokens of context, with $\eta_\beta = 0.8$, $\eta_\pi = 0.6$ and 4~GB per GPU
set aside. Which input is most likely to make the prediction wrong?

\textbf{Solution.} The model has 70.55~billion parameters: 12.08~billion in
attention, 56.37~billion in the FFNs and 2.10~billion in the two tables. A step
multiplies everything but the input table, 69.5~GB in FP8. The cache is
$2 \times 80 \times 8 \times 128 \times 1 = 160$~KiB per token, 671~MB per
4,096-token sequence. Eight GPUs are, to the napkin model, one machine with
26.8~TB/s, 15.8~PFLOP/s of dense FP8 and 640~GB. The batch that fits is
$(640 - 70.6 - 32) / 0.671 = 800$ sequences. At $B = 800$ a step reads
$69.5 + 800 \times 0.671 = 606$~GB, 28.3~ms at 21.4~TB/s, and computes
$800 \times 150$~GFLOP, 12.6~ms at 9.5~PFLOP/s: bandwidth-bound, 28,300 tokens per
second in aggregate, 35 per user, and at eight times \$3.99 an hour, \$0.31 per
million tokens --- the same price as Qwen3-8B on one GPU in BF16, because in both
deployments each GPU ends up holding about a hundred 4,096-token caches and
needs about 28~ms to stream its full memory. At $B = 1$ the step
reads 70.2~GB, 3.3~ms: 306 tokens per second for one user. The input most likely
to be wrong is the one the problem told us to ignore. Tensor parallelism needs two
all-reduces of the hidden states per layer, 160 per step. At batch one each moves
only 16~KB but costs a synchronization of eight GPUs, several microseconds each,
so they add a millisecond or more to a 3.3~ms step; at $B = 800$ each GPU moves
about 3.7~GB of partial sums per step, several milliseconds even at NVLink's
900~GB/s (\Chref{model-parallelism}).

\subsection{Goodput with two objectives}
\textbf{Problem.} A service runs Qwen3-8B in BF16 on one H100 and promises a
time to first token under 500~ms and a time per output token under 30~ms, for
requests of 2,000 prompt tokens and 300 output tokens. With the chapter's
efficiencies and chunked prefill that mixes prompt tokens into decode steps, how
many requests per second can it accept, and which objective binds first?

\textbf{Solution.} A request's prompt costs $2 \times 7.57 \times 10^{9} \times
2{,}000$ FLOPs of matrix products plus about 1.2~TFLOP of causal attention,
31.5~TFLOP in all: 53~ms of the GPU at 60\% of peak. During decode the mean context
is about 2,150 tokens, so a decode step with $B$ sequences takes
$5.65 + 0.118B$~ms of memory time (its compute is a quarter of that). At an
arrival rate $\lambda$, each step of length $T$ must also absorb
$\lambda T \times 53$~ms of prefill work, and by Little's law the number of
sequences decoding is $B = \lambda \times 300 \times T$. Setting $T = 30$~ms:
$30(1 - 0.053\lambda) = 5.65 + 0.118 \times 9\lambda$, so $\lambda = 9.2$ requests
per second, with $B = 82$ sequences decoding (derived). Memory is not the limit:
188 sequences of 2,150 tokens would fit. Nor is the first-token objective: at
that rate each 30~ms step carries about 550 prompt tokens, so a new prompt is
absorbed in about four steps, 120~ms plus queueing, far inside 500~ms. The time
per output token binds first, at about 2,750 output tokens and 18,300 prompt tokens
per second. Half of every step is prefill; a service with longer replies or
shorter prompts would get more requests through, and one with a tighter TPOT
objective would get fewer --- the trade \Chref{scheduling} makes explicit.

\subsection{DeepSeek's decode step, line by line}
\textbf{Problem.} DeepSeek reports that each decode GPU in its V3/R1 service
held two routed experts and the shared expert, produced about 1,850 tokens per
second while users saw 20--22, and that the average cache length per output
token was 4,989 \cite{deepseek2025inference}. From DeepSeek-V3's configuration
(61 layers, the first 3 dense; width 7,168; 128 heads with multi-head latent
attention caching 576 numbers per token per layer in 16 bits; experts of
$3 \times 7{,}168 \times 2{,}048$ weights, 8 routed per token; FP8 weights), price one
GPU's decode step in memory, compute and network time, and compare with the
reported step.

\textbf{Solution.} Throughput over per-user speed gives about 88 conversations
per GPU, stepping every $1/21 \approx 48$~ms if each step yields one token per
conversation (the report does not say whether multi-token prediction was used; if
it was, steps are further apart). \emph{Memory:} the GPU holds the full attention
weights for its own conversations, about 11.4~GB in FP8; the shared expert,
2.6~GB; the three dense layers' FFNs, 1.2~GB; its two routed experts in each of 58
layers, 5.1~GB; the output head, 0.9~GB --- 21~GB of weights per step. The cache
is $88 \times 4{,}989 \times 70{,}272$~bytes $= 31$~GB. The 52~GB take 15.5~ms at
3.35~TB/s, assuming the H800 streams memory as fast as the H100 SXM whose
datasheet gives that figure. \emph{Compute:} attention projections for 88 tokens,
2.0~TFLOP; the routed experts, which receive $88 \times 144 \times 8 / 288 = 352$
tokens per expert slot per layer, 3.6~TFLOP; shared expert, dense layers and output
head, 0.8~TFLOP --- 6.4~TFLOP in FP8, 3.2~ms at 1,979~TFLOP/s; and the latent
attention itself, 1.4~GFLOP per token per layer at 4,989 tokens of context,
7.5~TFLOP in BF16, 7.5~ms at 989~TFLOP/s --- 10.8~ms in all. \emph{Network:} each
token's hidden state goes to its 8 experts (7~KB each in FP8) and comes back
(14~KB in BF16) in each of 58 layers: $88 \times 8 \times 21{,}504 \times 58 =
0.88$~GB per step, 17.6~ms over a 50~GB/s InfiniBand link \cite{deepseek2024v3}
(an upper bound, since one expert in eighteen is on the same node). The floors
sum to 44~ms against about 48~ms reported. The step is not memory-bound, or
compute-bound, or network-bound: it uses all three, for a quarter to two fifths
of its length each, and it is only as short as it is because DeepSeek's five-stage
pipeline overlaps them --- incompletely, by this arithmetic (derived; every
input is the operator's report or the model's configuration, and nothing here
was measured by the book).
